\documentclass[10pt,a4paper]{amsart}
\usepackage[margin=19mm]{geometry}
\usepackage{amsmath,amssymb,mathtools}
\usepackage[scr=rsfs,cal=euler]{mathalfa}
\usepackage{xfrac}
\usepackage[unicode,hidelinks,bookmarks=false]{hyperref}
\newcommand{\ie}{i.e.\ }
\newcommand{\cf}{cf.\ }
\newcommand{\resp}{respectively\ }
\newcommand{\Subsection}{\subsection}
\providecommand{\nobreakdash}{\nobreak\hbox{-}}
\setlength{\emergencystretch}{2em}
\multlinegap0pt
\usepackage{bm}
\usepackage{amssymb}
\usepackage{dsfont}
\usepackage{mathtools}
\usepackage{tikz}
\newtheorem{theorem}{Theorem}
\newcommand{\A}{\mathrm{Aut}}
\newcommand{\Aa}{\A_{(I)}^{\M}(\K)}
\newcommand{\G}{\mathrm{G}}
\newcommand{\K}{\mathcal{K}}
\newcommand{\M}{\mathcal{M}}
\newcommand{\pa}{\mathrm{PA}}
\newcommand{\tp}{\mathrm{tp}}
\title{The automorphism group of countable recursively saturated models of Peano arithmetic and strong cuts}
\author{Saeideh Bahrami}
\date{Source: arXiv:2604.10282v1 (CC BY 4.0)}
\begin{document}
\maketitle

\noindent\emph{Source and licence.} The automorphism group of countable recursively saturated models of Peano arithmetic and strong cuts. Version arXiv:2604.10282v1 (CC BY 4.0), distributed by arXiv under the Creative Commons Attribution 4.0 International licence, http://creativecommons.org/licenses/by/4.0/, as recorded in the arXiv licence metadata for this exact version and shown on the abstract page https://arxiv.org/abs/2604.10282v1. Full licence notice: \texttt{licenses/arxiv-2604.10282-CC-BY-4.0.txt}.\par\smallskip

\noindent\textit{Editorial scope.} Counted unit: the article's theorem theorem (source0.tex lines 305--307). The article's complete original proof of it is delivered with it (source0.tex lines 308--402, one \texttt{proof} environment of 95 lines). No mathematics is written, repaired or re-derived: the statement and the proof are the article's own lines, in the article's own notation, and no retained line is substituted or re-flowed.  No further source-local context is needed: every cross-reference of every retained line resolves inside the two delivered spans.  Nothing was omitted from any retained span, so the package carries no omission record.  No retained line contains a citation, so the wrapper carries no bibliography and no bibliography entry.  No retained line contains a figure environment, an image-inclusion command or drawing code, so no image is delivered and the package declares no asset dependency.  Editorial formatting only, in addition: an AMS \texttt{amsart} preamble that restates the article's own declarations and macros used by the retained lines, namely the environments \texttt{theorem} on separate counters, together with the standard packages those lines need.  The retained environments print in this document as Theorem 1 in order of appearance, whereas the article prints the same items with the numbering of its own full document.  The complete native source of the article as distributed in the arXiv e-print for this exact version is bundled unchanged as \texttt{sources/198258/source0.tex}.
\begin{theorem}
	($n\in\omega$)	Suppose $ \M $ is a countable  recursively saturated model of $ \pa $, $I$ is a strong cut of $\M$, and $ a,b_{0},...,b_{n}\in M $ such that $ \tp(a,i)=\tp(b_{k},i) $ for every $k=0,...,n$ and for all $i\in I$. Then there exists some $I$-small elementary submodel $ \K $ of $\M$  and some $I$-e.c.  tuple $ (g_{0},...,g_{n})\in(\Aa)^{n+1} $ such that $ a,b_{0},...,b_{n}\in K $ and $ g_{k}(a)=b_{k} $ for every $k=0,...,n$.
\end{theorem}

\begin{proof}
	We will prove the theorem for the case $n=0$, and the remaining cases follow by a similar argument.  So for simplicity, let $b:=b_{0}$.\\
		Let $S$ be an inductive satisfaction class for $\M$ such that $(\M;S)$ is also recursively saturated. So by Lemma 2.2(1), there exists some  $I$-small elementary submodel  $(\K;S')$  of $(\M;S)$ containing $a,b$.  Let $\K=\{(\alpha)_{i}: \ i\in I\}$ for some $\alpha\in M$. Moreover, let $a=(\alpha)_{i_{a}}$ and $b=(\alpha)_{i_{b}}$ for some $i_{a},i_{b}\in I$. Now, we  define:
	$$\Delta(x,y):=\mathrm{max}\{z: \ \forall \langle w,i\rangle<z \ S(\bm{\phi}_{w}(x,i))\leftrightarrow S(\bm{\phi}_{w}(y,i))\}\footnote{Throughout this paper, we adopt the convention that $\max(\emptyset)=0$.};$$ 
	and
	$$\Lambda(\langle i,j\rangle):=\Delta((\alpha)_{i},(\alpha)_{j})\footnote{For case $n>0$, we define: $$\Lambda(\langle i_{0},i_{1},...,i_{n+1}\rangle):=\min\{\Delta((\alpha)_{i_{k}},(\alpha)_{i_{m}}): \ 0\leq k\leq m\leq n+1  \}. $$}. $$
	
	Since $I$ is strong in $\M$,	there exists some $\epsilon>I$ such that $\Lambda(i)\in I$ iff $\Lambda(i)<\epsilon $ for all $i\in I$.  Then, let	$\langle\varphi_{r}((\alpha)_{i},x,y): \  r,i\in M\rangle$ be a  canonical enumeration of $\mathbb{L}_{A}$-formulas of the given form with $(\alpha)_{i}$s as parameters.  Moreover, let $\Theta(i,u,v,x,y)$ be the following $\mathbb{L}$-formula (where $\mathbb{L}:=\mathbb{L}_{A}\cup\{S\}$):
	$$  \left(\begin{array}{c}
		S(\varphi_{(i)_{0}}((\alpha)_{(i)_{1}},(\alpha)_{x},(\alpha)_{y})) \ \wedge \\ \forall z,w  \left(\begin{array}{c}
	   	(\alpha)_{z}=(\alpha)_{u}^{\frown}(\alpha)_{x}\ \wedge  \
		(\alpha)_{w}=(\alpha)_{v}^{\frown}(\alpha)_{y}
	\end{array}\right)  
	\rightarrow\Lambda(\langle z,w\rangle)>\epsilon\end{array}\right). $$
	Now,  we inductively define  $\mathbb{L}$-terms $\Omega(i)$ and $\bar{\Omega}(i)$ in $\M$ as follows:
	\begin{center}
	$\Omega(0):=\langle i_{a},i_{b}\rangle$, and\\ $\bar{\Omega}(0):=\langle \textbf{u},\textbf{v}\rangle\in I$ such that $(\alpha)_{\textbf{u}}=\langle a\rangle$ and $(\alpha)_{\textbf{v}}=\langle b\rangle$.
	\end{center} 
	and
	\begin{equation*}
		\Omega(i+1):=  \left\{
		\begin{array}{rl}
			\mu_{\langle x,y\rangle}   \Theta( i,(\bar{\Omega}(i))_{0},(\bar{\Omega}(i))_{1},x,y) & \text{if } \scriptstyle(\M;S)\models\exists x,y \    \Theta( i,(\bar{\Omega}(i))_{0},(\bar{\Omega}(i))_{1},x,y)\\
			\langle i_{a},i_{b}\rangle \   \ \ \ \ \ \ \ \ \ \ \ \ \ \ \ \ \ \ \ \ \ \ \ \ \ \ \ \ \ \ \    & \text{o.w.} 
		\end{array} \right.
	\end{equation*}
\begin{equation*}
	\bar{\Omega}(i+1):= \left\{
	\begin{array}{rl}
	\mu_{\langle u,v \rangle}  \left(\begin{array}{c}
		(\alpha)_{u}=((\alpha)_{(\bar{\Omega}(i))_{0}})^{\frown}(\alpha)_{(\Omega(i+1))_{0}}  \ \wedge \\ (\alpha)_{v}=((\alpha)_{(\bar{\Omega}(i))_{1}})^{\frown}(\alpha)_{(\Omega(i+1))_{1}}
	\end{array}\right) & \text{if }\scriptsize\exists u,v 
 \left(\begin{array}{c}
	(\alpha)_{u}=((\alpha)_{(\bar{\Omega}(i))_{0}})^{\frown}(\alpha)_{(\Omega(i+1))_{0}}  \ \wedge \\ (\alpha)_{v}=((\alpha)_{(\bar{\Omega}(i))_{1}})^{\frown}(\alpha)_{(\Omega(i+1))_{1}}
\end{array}\right)\\
\epsilon \ \ \ \ \ \ \ \ \ \ \ \ \ \ \ \ \ \ \ \ \ \ \ \ \ \ \ \ \ \ \ \ \ \ \ \ \ \ \ \ \ \ \ \ \ \ \ \ \ \ \ \ \ \ \ \ \  & \text{o.w.}
\end{array} \right.
\end{equation*}	
	Again, since $I$ is strong in $\M$	there exists some $\eta,\varepsilon>I$ such that $\Omega(i)\in I$ iff ${\Omega(i)<\varepsilon} $, and  $\bar{\Omega}(i)\in I$ iff $\bar{\Omega}(i)<\eta $ for all $i\in I$. Note that $\Omega(i)<\varepsilon $ and $\bar{\Omega}(i)<\eta$ for all $i\in I$: we use induction to show this fact. Suppose we have shown $\Omega(i)<\varepsilon $ and $\bar{\Omega}(i)<\eta$ for $i\in I$. By Lemma 2.2(3), it suffices to show that $\Omega(i+1)<\varepsilon$. Suppose ${(\M;S)\models\exists x,y \    \Theta( i,(\bar{\Omega}(i))_{0},(\bar{\Omega}(i))_{1},x,y)}$ (for the other case, it is clear that $\Omega(i+1)<\varepsilon$). So, for every $s\in I$ it holds that:\\
	
$(1): \ \ \ $	$(\M;S)\models\exists x,y  \left(\begin{array}{c}
		S(\varphi_{(i)_{0}}((\alpha)_{(i)_{1}},x,y)) \ \wedge \\ 
		\forall u,v  \left(
		\begin{array}{c}
			u=((\alpha)_{(\bar{\Omega}(i))_{0}})^{\frown}x \ \wedge \\ v=((\alpha)_{(\bar{\Omega}(i))_{1}})^{\frown}y	\end{array}\right)\rightarrow
			\forall \langle r,z\rangle<s \
		 S(\bm{\phi}_{r}(u,z))\leftrightarrow S(\bm{\phi}_{r}(v,z)) 
		\end{array}\right)$.\\
	
	So since $(\K;S')\prec(\M;S)$, statement (1) implies that for every $s\in I$ we have:\\
	
	$(2): \ \ \ $	$(\K;S')\models\exists x,y  \left(\begin{array}{c}
		S'(\varphi_{(i)_{0}}((\alpha)_{(i)_{1}},x,y)) \ \wedge \\ 
		\forall u,v \left(
		\begin{array}{c}
			u=((\alpha)_{(\bar{\Omega}(i))_{0}})^{\frown}x \ \wedge \\ v=((\alpha)_{(\bar{\Omega}(i))_{1}})^{\frown}y	\end{array}\right)\rightarrow
		\forall \langle r,z\rangle<s \
		S'(\bm{\phi}_{r}(u,z))\leftrightarrow S'(\bm{\phi}_{r}(v,z)) 
	\end{array}\right)$.\\
	
As a result, by using  Overspill principle over $ I$ in $(\K;S')$ and Lemma 2.2(3), we infer that $\Omega(i+1)\in I$.	
	
	
		Now, we define the following recursive type for every $s\in M$:
	\begin{align*}
		p_{s}(x,y):=&\{\forall i<s \ \phi(x,i)\leftrightarrow\phi(y,i): \ \phi \text{ is an } \mathbb{L}_{A}\text{-formula}\}\cup\\&
	\left\lbrace	
		\forall i<s \ 
	 ((x)_{i}=(\alpha)_{(\Omega(i))_{0}} \ \wedge\ (y)_{i}=(\alpha)_{(\Omega(i))_{1}}) 
	\right\rbrace.
\end{align*}	
	
	
	First, we will prove that there exists some  $s>I$ such that  $p_{s}(x,y)$ is finitely satisfiable. For this purpose, define:
	\begin{center}
		$\Psi(m):= \max\left\lbrace s: \exists x,y \ 	\psi(x,y,s,m,\alpha) \right\rbrace$, where $\psi(x,y,s,m,\alpha) $ is the following $\mathbb{L}$-formula:
		$$\forall i<s \ \forall z<m \left(\begin{array}{c}
			(S(\bm{\phi}_{z}(x,i))\leftrightarrow S(\bm{\phi}_{z}(y,i))) \ \wedge\\
				 ((x)_{i}=(\alpha)_{(\Omega(i))_{0}} \ \wedge\ (y)_{i}=(\alpha)_{(\Omega(i))_{1}}) 
		\end{array}\right)  .$$
	\end{center}
	
	Let $e>I$ be the witness of strength of $I$ in $\M$ for the definable function $\Psi$. We will  show that $p_{e}(x,y)$ is finitely satisfiable. For this purpose, it suffices to prove that $\Psi(m)>e$ for every $m\in\omega$. So, suppose $m\in\omega$ be arbitrary. By the way we defined $\Omega(i)$ and by using induction in $(\M;S)$, we can show  that ${(\M;S)\models\exists x,y \ \psi(x,y,s,m,\alpha)}$ for every $s\in I$. So, by Overspill principle over $I$ in $(\M;S)$, we have $\Psi(m)>e$.
	
	As a result,  we can find some $ \lambda,\xi\in M $ such that $(\lambda,\xi)$ realizes $p_{e}(x,y)$. Now, put:
	$$ g:=\bigcup_{i\in I}(\lambda)_{i}\mapsto(\xi)_{i} .$$ 
	Then, we have the following claim:
	\begin{center}
		\textit{Claim:} For every $m\in\omega$ and for each $i\in I$, if  $\M\models\exists x \ \varphi_{m}((\alpha)_{i},x,f(x))$ for some $f\in\G_{I}$ such that $g\subseteq f$, then $\M\models\varphi_{m}((\alpha)_{i},(\lambda)_{j},(\xi)_{j})$ for some $j\in I$.
	\end{center}
	
The claim implies that, first $\mathrm{Dom}(g)=K$:  note that since $\Omega(i)\in I$ for all $i\in I$, it holds that $\mathrm{Dom}(g)\subseteq K $. To see $K\subseteq \mathrm{Dom}(g)$, let $i\in I$ be arbitrary and consider  the $\mathbb{L}_{A}$-formula $(x=(\alpha)_{i} \ \wedge \ y=y)$ and $f\in\G_{(I)} $ such that $f(\lambda)=\xi$ in the claim. So we conclude that $(\alpha)_{i}\in \mathrm{Dom}(g)$. In a similar way, $\mathrm{Range}(g)=K$. Secondly, the claim implies that $g\in\Aa$ is $I$-e.c.
	
	In order to prove the claim, suppose we have the assumptions of the Claim. Thus, by iterating similar arguments to statements (1) and (2), we conclude that $\M\models\varphi_{m}((\alpha)_{i},(\lambda)_{j},(\xi)_{j})$ where $j:=\Omega(\langle m,i\rangle+1)$.
\end{proof}

\end{document}
